% --- Archon physics formalization source begin ---
% archon:physics
% archon:covers PhyXMiniProblems/problem_phyx_mini_0769.lean
% archon:source-report reports/phyx_mini/problem_phyx_mini_0769.source.json

\chapter{Physics problem phyx\_mini\_0769}
\label{ch:PhyXMiniProblems_problem_phyx_mini_0769}

\paragraph{Problem source.}
\#\# Physical scenario

A 45.0 kg woman stands up in a 60.0 kg canoe 5.00 m long. She walks from a point 1.00 m from one end to a point 1.00 m from the other end as shown in figure.

\#\# Question

If you ignore resistance to motion of the canoe in the water, how far does the canoe move during this process?

\#\# Answer choices

- A: A: 1.87m

- B: B: 1.29m

- C: C: 0.12m

- D: D: 0.98m

\#\# Figure caption (auxiliary; use the image as primary evidence)

The image depicts a woman walking inside a canoe on water. The canoe is oriented horizontally, and the woman is moving from one end to the other. 

Below the canoe, there is a horizontal arrow indicating movement from left to right. The labeled sections are:
- "Start" to the left, marked with a distance of "1.00 m."
- A central section of "3.00 m."
- "Finish" to the right, also marked with a distance of "1.00 m."

The background shows light ripples, suggesting the canoe is on water. There are no additional objects or text present.

\#\# Dataset metadata

Momentum and Collisions; Physical Model Grounding Reasoning, Spatial Relation Reasoning

\paragraph{Recorded answer/context.}
B: B: 1.29m

\paragraph{Figure/image path.}
/root/proposal\_for\_physic/science-mango/phyx\_mini\_run/phyx\_data/test\_image/769.png

\paragraph{Formalization target.}
create a compiling Lean file with sorry bodies at `PhyXMiniProblems/problem\_phyx\_mini\_0769.lean`. The Lean declarations must preserve the physical quantities, dimensions or dimensional roles, figure labels, governing-law hypotheses, and final relation expressed by this problem.
Use Mathlib/Physlib names found through LeanExplore where available. If a physics API is missing, introduce faithful local abstractions rather than scalar placeholder aliases.

\begin{theorem}[Physics formalization target]
\label{thm:physics:phyx_mini_0769:target}
\lean{PhyXMiniProblems.ProblemPhyXMini0769.problem_phyx_mini_0769}
\uses{def:physics:phyx-mini-0769:phyxminiproblems-problemphyxmini0769-displacementinmeters, def:physics:phyx-mini-0769:phyxminiproblems-problemphyxmini0769-canoewalksetup, def:physics:phyx-mini-0769:phyxminiproblems-problemphyxmini0769-matchescanoewalkscenario, def:physics:phyx-mini-0769:phyxminiproblems-problemphyxmini0769-matchesproblemreadouts, def:physics:phyx-mini-0769:phyxminiproblems-problemphyxmini0769-matchesprimaryfigure, def:physics:phyx-mini-0769:phyxminiproblems-problemphyxmini0769-satisfiesrigidcanoekinematics, def:physics:phyx-mini-0769:phyxminiproblems-problemphyxmini0769-satisfiesisolatedhorizontalcenterofmasslaw, def:physics:phyx-mini-0769:phyxminiproblems-problemphyxmini0769-displayeddistanceinmeters, def:physics:phyx-mini-0769:phyxminiproblems-problemphyxmini0769-canoetraveldistanceinmeters, def:physics:phyx-mini-0769:phyxminiproblems-problemphyxmini0769-agreeswithdisplayedhundredth, def:physics:phyx-mini-0769:phyxminiproblems-problemphyxmini0769-isuniqueclosestdistancechoice}
The assigned autoformalize agent should translate this physics problem into Lean declarations in the covered file, with theorem and lemma proof bodies written as `by sorry`.
\end{theorem}
\begin{proof}
This is an autoformalization task, not a proof task. Produce statements that can later be proved by the physics prover without weakening the physical model.
\end{proof}

% --- Archon declaration topology begin ---
\section{Declaration topology}
\begin{definition}[Mass Quantity]
  \label{def:physics:phyx-mini-0769:phyxminiproblems-problemphyxmini0769-massquantity}
  \lean{PhyXMiniProblems.ProblemPhyXMini0769.MassQuantity}
  A unit-independent nonnegative physical mass
\end{definition}
\begin{proof}
  This is the declared model definition of Mass Quantity.
\end{proof}
\begin{definition}[Length Magnitude]
  \label{def:physics:phyx-mini-0769:phyxminiproblems-problemphyxmini0769-lengthmagnitude}
  \lean{PhyXMiniProblems.ProblemPhyXMini0769.LengthMagnitude}
  A unit-independent nonnegative physical length
\end{definition}
\begin{proof}
  This is the declared model definition of Length Magnitude.
\end{proof}
\begin{definition}[Horizontal Displacement Quantity]
  \label{def:physics:phyx-mini-0769:phyxminiproblems-problemphyxmini0769-horizontaldisplacementquantity}
  \lean{PhyXMiniProblems.ProblemPhyXMini0769.HorizontalDisplacementQuantity}
  A signed one-dimensional displacement along the canoe's horizontal axis
\end{definition}
\begin{proof}
  This is the declared model definition of Horizontal Displacement Quantity.
\end{proof}
\begin{definition}[mass In Kilograms]
  \label{def:physics:phyx-mini-0769:phyxminiproblems-problemphyxmini0769-massinkilograms}
  \lean{PhyXMiniProblems.ProblemPhyXMini0769.massInKilograms}
  \uses{def:physics:phyx-mini-0769:phyxminiproblems-problemphyxmini0769-massquantity}
  Coherent-SI kilogram readout of a physical mass
\end{definition}
\begin{proof}
  This is the declared model definition of mass In Kilograms.
\end{proof}
\begin{definition}[length In Meters]
  \label{def:physics:phyx-mini-0769:phyxminiproblems-problemphyxmini0769-lengthinmeters}
  \lean{PhyXMiniProblems.ProblemPhyXMini0769.lengthInMeters}
  \uses{def:physics:phyx-mini-0769:phyxminiproblems-problemphyxmini0769-lengthmagnitude}
  Coherent-SI metre readout of a nonnegative physical length
\end{definition}
\begin{proof}
  This is the declared model definition of length In Meters.
\end{proof}
\begin{definition}[displacement In Meters]
  \label{def:physics:phyx-mini-0769:phyxminiproblems-problemphyxmini0769-displacementinmeters}
  \lean{PhyXMiniProblems.ProblemPhyXMini0769.displacementInMeters}
  \uses{def:physics:phyx-mini-0769:phyxminiproblems-problemphyxmini0769-horizontaldisplacementquantity}
  Signed coherent-SI metre readout along the image's horizontal axis
\end{definition}
\begin{proof}
  This is the declared model definition of displacement In Meters.
\end{proof}
\begin{definition}[Canoe Station]
  \label{def:physics:phyx-mini-0769:phyxminiproblems-problemphyxmini0769-canoestation}
  \lean{PhyXMiniProblems.ProblemPhyXMini0769.CanoeStation}
  The four geometrically distinguished stations along the canoe
\end{definition}
\begin{proof}
  This is the declared model definition of Canoe Station.
\end{proof}
\begin{definition}[Figure Segment]
  \label{def:physics:phyx-mini-0769:phyxminiproblems-problemphyxmini0769-figuresegment}
  \lean{PhyXMiniProblems.ProblemPhyXMini0769.FigureSegment}
  The three successive length annotations under the canoe
\end{definition}
\begin{proof}
  This is the declared model definition of Figure Segment.
\end{proof}
\begin{definition}[endpoints]
  \label{def:physics:phyx-mini-0769:phyxminiproblems-problemphyxmini0769-figuresegment-endpoints}
  \lean{PhyXMiniProblems.ProblemPhyXMini0769.FigureSegment.endpoints}
  \uses{def:physics:phyx-mini-0769:phyxminiproblems-problemphyxmini0769-canoestation, def:physics:phyx-mini-0769:phyxminiproblems-problemphyxmini0769-figuresegment}
  Endpoints associated with each of the three displayed dimension segments
\end{definition}
\begin{proof}
  This is the declared model definition of endpoints.
\end{proof}
\begin{definition}[Horizontal Direction]
  \label{def:physics:phyx-mini-0769:phyxminiproblems-problemphyxmini0769-horizontaldirection}
  \lean{PhyXMiniProblems.ProblemPhyXMini0769.HorizontalDirection}
  Qualitative horizontal directions in the supplied image
\end{definition}
\begin{proof}
  This is the declared model definition of Horizontal Direction.
\end{proof}
\begin{definition}[station Has Printed Text]
  \label{def:physics:phyx-mini-0769:phyxminiproblems-problemphyxmini0769-stationhasprintedtext}
  \lean{PhyXMiniProblems.ProblemPhyXMini0769.stationHasPrintedText}
  \uses{def:physics:phyx-mini-0769:phyxminiproblems-problemphyxmini0769-canoestation}
  Exactly the two interior stations have printed text labels in the image
\end{definition}
\begin{proof}
  This is the declared model definition of station Has Printed Text.
\end{proof}
\begin{definition}[Canoe Walk Figure]
  \label{def:physics:phyx-mini-0769:phyxminiproblems-problemphyxmini0769-canoewalkfigure}
  \lean{PhyXMiniProblems.ProblemPhyXMini0769.CanoeWalkFigure}
  \uses{def:physics:phyx-mini-0769:phyxminiproblems-problemphyxmini0769-canoestation, def:physics:phyx-mini-0769:phyxminiproblems-problemphyxmini0769-figuresegment}
  Literal and qualitative content transcribed from image 769.png
\end{definition}
\begin{proof}
  This is the declared model definition of Canoe Walk Figure.
\end{proof}
\begin{definition}[Canoe Walk Setup]
  \label{def:physics:phyx-mini-0769:phyxminiproblems-problemphyxmini0769-canoewalksetup}
  \lean{PhyXMiniProblems.ProblemPhyXMini0769.CanoeWalkSetup}
  \uses{def:physics:phyx-mini-0769:phyxminiproblems-problemphyxmini0769-massquantity, def:physics:phyx-mini-0769:phyxminiproblems-problemphyxmini0769-lengthmagnitude, def:physics:phyx-mini-0769:phyxminiproblems-problemphyxmini0769-horizontaldisplacementquantity, def:physics:phyx-mini-0769:phyxminiproblems-problemphyxmini0769-canoestation, def:physics:phyx-mini-0769:phyxminiproblems-problemphyxmini0769-figuresegment, def:physics:phyx-mini-0769:phyxminiproblems-problemphyxmini0769-horizontaldirection, def:physics:phyx-mini-0769:phyxminiproblems-problemphyxmini0769-canoewalkfigure}
  Independent physical quantities and scenario roles. In particular, neither horizontal displacement is defined from the numerical answer; the governing relations below constrain both displacement fields
\end{definition}
\begin{proof}
  This is the declared model definition of Canoe Walk Setup.
\end{proof}
\begin{definition}[Matches Canoe Walk Scenario]
  \label{def:physics:phyx-mini-0769:phyxminiproblems-problemphyxmini0769-matchescanoewalkscenario}
  \lean{PhyXMiniProblems.ProblemPhyXMini0769.MatchesCanoeWalkScenario}
  \uses{def:physics:phyx-mini-0769:phyxminiproblems-problemphyxmini0769-canoewalksetup}
  Qualitative scenario conditions stated or physically implicit in the problem
\end{definition}
\begin{proof}
  This is the declared model definition of Matches Canoe Walk Scenario.
\end{proof}
\begin{definition}[Matches Problem Readouts]
  \label{def:physics:phyx-mini-0769:phyxminiproblems-problemphyxmini0769-matchesproblemreadouts}
  \lean{PhyXMiniProblems.ProblemPhyXMini0769.MatchesProblemReadouts}
  \uses{def:physics:phyx-mini-0769:phyxminiproblems-problemphyxmini0769-massinkilograms, def:physics:phyx-mini-0769:phyxminiproblems-problemphyxmini0769-lengthinmeters, def:physics:phyx-mini-0769:phyxminiproblems-problemphyxmini0769-canoewalksetup}
  Numerical mass and total-length readouts from the problem prose
\end{definition}
\begin{proof}
  This is the declared model definition of Matches Problem Readouts.
\end{proof}
\begin{definition}[Matches Primary Figure]
  \label{def:physics:phyx-mini-0769:phyxminiproblems-problemphyxmini0769-matchesprimaryfigure}
  \lean{PhyXMiniProblems.ProblemPhyXMini0769.MatchesPrimaryFigure}
  \uses{def:physics:phyx-mini-0769:phyxminiproblems-problemphyxmini0769-lengthinmeters, def:physics:phyx-mini-0769:phyxminiproblems-problemphyxmini0769-stationhasprintedtext, def:physics:phyx-mini-0769:phyxminiproblems-problemphyxmini0769-canoewalksetup}
  Primary-image evidence. The dimension line is a measurement line, not a claim about the canoe's direction of motion. No canoe displacement or answer choice occurs in this record
\end{definition}
\begin{proof}
  This is the declared model definition of Matches Primary Figure.
\end{proof}
\begin{definition}[Satisfies Rigid Canoe Kinematics]
  \label{def:physics:phyx-mini-0769:phyxminiproblems-problemphyxmini0769-satisfiesrigidcanoekinematics}
  \lean{PhyXMiniProblems.ProblemPhyXMini0769.SatisfiesRigidCanoeKinematics}
  \uses{def:physics:phyx-mini-0769:phyxminiproblems-problemphyxmini0769-lengthinmeters, def:physics:phyx-mini-0769:phyxminiproblems-problemphyxmini0769-displacementinmeters, def:physics:phyx-mini-0769:phyxminiproblems-problemphyxmini0769-canoewalksetup}
  Rigid-body translation of the canoe makes the woman's ground displacement equal to the canoe displacement plus her displacement relative to the canoe. This is a kinematic law, not the requested canoe displacement formula
\end{definition}
\begin{proof}
  This is the declared model definition of Satisfies Rigid Canoe Kinematics.
\end{proof}
\begin{definition}[Satisfies Isolated Horizontal Center Of Mass Law]
  \label{def:physics:phyx-mini-0769:phyxminiproblems-problemphyxmini0769-satisfiesisolatedhorizontalcenterofmasslaw}
  \lean{PhyXMiniProblems.ProblemPhyXMini0769.SatisfiesIsolatedHorizontalCenterOfMassLaw}
  \uses{def:physics:phyx-mini-0769:phyxminiproblems-problemphyxmini0769-massinkilograms, def:physics:phyx-mini-0769:phyxminiproblems-problemphyxmini0769-displacementinmeters, def:physics:phyx-mini-0769:phyxminiproblems-problemphyxmini0769-canoewalksetup}
  For the initially resting woman--canoe system with no external horizontal impulse, conservation of horizontal momentum leaves the center of mass fixed. The resulting first-moment balance is stated directly because the searched rigid-body APIs do not provide a two-body, finite-walk interface with these dimensionful observables
\end{definition}
\begin{proof}
  This is the declared model definition of Satisfies Isolated Horizontal Center Of Mass Law.
\end{proof}
\begin{definition}[Answer Choice]
  \label{def:physics:phyx-mini-0769:phyxminiproblems-problemphyxmini0769-answerchoice}
  \lean{PhyXMiniProblems.ProblemPhyXMini0769.AnswerChoice}
  The four distance values printed in the multiple-choice list, in metres
\end{definition}
\begin{proof}
  This is the declared model definition of Answer Choice.
\end{proof}
\begin{definition}[displayed Distance In Meters]
  \label{def:physics:phyx-mini-0769:phyxminiproblems-problemphyxmini0769-displayeddistanceinmeters}
  \lean{PhyXMiniProblems.ProblemPhyXMini0769.displayedDistanceInMeters}
  \uses{def:physics:phyx-mini-0769:phyxminiproblems-problemphyxmini0769-answerchoice}
  Displayed distance in metres associated with each supplied answer choice
\end{definition}
\begin{proof}
  This is the declared model definition of displayed Distance In Meters.
\end{proof}
\begin{definition}[canoe Travel Distance In Meters]
  \label{def:physics:phyx-mini-0769:phyxminiproblems-problemphyxmini0769-canoetraveldistanceinmeters}
  \lean{PhyXMiniProblems.ProblemPhyXMini0769.canoeTravelDistanceInMeters}
  \uses{def:physics:phyx-mini-0769:phyxminiproblems-problemphyxmini0769-displacementinmeters, def:physics:phyx-mini-0769:phyxminiproblems-problemphyxmini0769-canoewalksetup}
  Distance traveled by the canoe, obtained from its signed displacement
\end{definition}
\begin{proof}
  This is the declared model definition of canoe Travel Distance In Meters.
\end{proof}
\begin{definition}[Agrees With Displayed Hundredth]
  \label{def:physics:phyx-mini-0769:phyxminiproblems-problemphyxmini0769-agreeswithdisplayedhundredth}
  \lean{PhyXMiniProblems.ProblemPhyXMini0769.AgreesWithDisplayedHundredth}
  Agreement of an exact metre value with a value displayed to hundredths
\end{definition}
\begin{proof}
  This is the declared model definition of Agrees With Displayed Hundredth.
\end{proof}
\begin{definition}[Is Unique Closest Distance Choice]
  \label{def:physics:phyx-mini-0769:phyxminiproblems-problemphyxmini0769-isuniqueclosestdistancechoice}
  \lean{PhyXMiniProblems.ProblemPhyXMini0769.IsUniqueClosestDistanceChoice}
  \uses{def:physics:phyx-mini-0769:phyxminiproblems-problemphyxmini0769-answerchoice, def:physics:phyx-mini-0769:phyxminiproblems-problemphyxmini0769-displayeddistanceinmeters}
  A choice is uniquely closest to the exact distance among the four options
\end{definition}
\begin{proof}
  This is the declared model definition of Is Unique Closest Distance Choice.
\end{proof}
% --- Archon declaration topology end ---
% --- Archon physics formalization source end ---
